\documentclass[11pt]{article}
\usepackage{amsmath,amssymb,amsthm,mathtools,hyperref}

\begin{document}
\section*{Golgi cargo-transport mechanism}
\label{prob:010}
\label{rank:211}
\par\smallskip

\noindent\textbf{Domain.} Secretory trafficking. \quad \textbf{Retained status/source context.} Open quantitative mechanism-integration problem. Cisternal maturation has direct experimental support in yeast; the question is the relative contributions and transfer across cargoes, not whether maturation exists.
\paragraph{Definitions and assumptions.} Let $c_i(t)$ be tagged cargo amount in Golgi cisterna $i$, $a_i(t)$ its measured maturation marker state, and $j_{ij}(t)$ the nonnegative vesicular or tubular fractional transfer rate from cisterna $i$ to cisterna $j$, in $\mathrm{s}^{-1}$, for $i\ne j$; set $j_{ii}=0$. Cisternal birth, age progression and loss are specified stochastic events. The conserved tracer has negligible synthesis during a pulse-chase and calibrated bleaching. A model $M$ combines cargo transfer, cisternal maturation and optional retrograde recycling; $h_M$ maps states to spatial imaging signals. Use at least two cargo sizes and a trafficking perturbation that does not abolish measurement.
\paragraph{Problem statement.} Estimate the fraction of forward cargo displacement caused by each mechanism and identify which mixtures are distinguishable:
\[
\dot c_i=\sum_j(j_{ji}c_j-j_{ij}c_i),\qquad
\theta\sim_M\theta'\ \Longleftrightarrow\ P^u_{M,\theta}=P^u_{M,\theta'}
\quad\forall u\in\mathcal U.
\]
Age-dependent cisternal motion is included separately in $h_M$. In words, determine how a cargo moves while the compartments themselves change.
\paragraph{Observation-law notation.} Here $\mathcal U$ is the declared admissible observation/intervention design family and $\theta$ ranges over the declared parameter space of the fixed candidate model $M$. Let $X$ denote its full latent state process, including cargo, cisternal ages and positions, and any latent maturation variables. The measured marker state $a_i$ is an observation, not an identification of latent maturation state with cisternal age. Each instance specifies the marker observation rule and the observation kernel, whose spatial signal map is $h_M$. The measure $P^u_{M,\theta}$ is the induced joint law of all recorded cargo and marker observations under design $u$, on the declared observation space and horizon or asymptotic regime, after integrating over latent initial states, stochastic dynamics and measurement noise. Equality in the display is equality of these laws. The relation shown compares parameters within $M$; any comparison between distinct candidate models must likewise compare their induced laws on a common observation space under the same admissible designs.
\paragraph{Additional binding assumptions and resolution criterion.} The balance equation applies between cisternal birth, loss and fusion events; specify mass-conserving jump maps and track cargo exiting the modeled Golgi. Define each mechanism's share from labeled positive forward transport events or from a declared counterfactual attribution rule. Net displacement can vanish or contain cancellations, so its ratio is not automatically a well-defined mechanistic fraction. Observational equivalence concerns complete measurement laws, including noise.
\paragraph{Remaining instance-level definition choices.} The meaning of forward progress remains to be declared: its coordinate, orientation, units and reference frame, cargo weighting, aggregation and treatment of progress within a cargo-retaining cisterna. Spatial displacement, biochemical progression and topological progression are not interchangeable. The choice between positive-event accounting and counterfactual attribution also remains open. An event-accounting instance must declare its mechanism labels, prevent double counting and specify the normalization, including the zero-total convention. A counterfactual instance must declare its causal model, intervention versions and interaction-allocation rule, including whether contributions can be negative or nonadditive; a fitted mixture alone does not establish causal fractions. The interpretation and units of $c_i$ must also be declared. The displayed differential balance does not by itself specify a realized molecule-count jump process or a closed equation for unconditional expected cargo under random rates. Complete the evolving cisternal index set and age, position and maturation dynamics, and account for any labeled influx, newborn-cisterna contents and external tracer reservoirs as well as exit. Negligible synthesis does not imply zero influx of previously labeled cargo. No particular choices among these remaining conventions are fixed here.

\paragraph{Alternative formulations.} (i) In dispersed yeast Golgi, track individual cisternae and compare cargo residence with maturation-marker exchange. (ii) In stacked mammalian Golgi, jointly model diffusion through tubules, vesicle exchange and cisternal progression from super-resolved pulse-chase data. The geometry and identifiable parameters differ, so results from the first system are not automatically equivalent.
\paragraph{Obstructions and partial progress.} Live yeast imaging supports maturation. Marker turnover alone need not reveal cargo transfer, and correlated compartment movement can mimic directed transport. A static graph misses compartment birth and loss. The unresolved generalization is a validated mixture of mechanisms under specified conditions; absence of a unique decomposition should be reported as an inference limit.
\paragraph{Sources (supplied provenance; not all rechecked).}
\begin{itemize}
\item Wikipedia, List of unsolved problems in biology (supplied background only; not used as scientific verification). \url{https://en.wikipedia.org/wiki/List_of_unsolved_problems_in_biology}
\item Losev et al., Golgi maturation visualized in living yeast (2006; primary evidence for a solved component). \url{https://www.nature.com/articles/nature04717}
\end{itemize}


\section*{Binding catalogue conventions}
This edition revises both supplied catalogues and preserves original problem numbers 1--330. It adds one globally unique label \verb|\label{rank:x}| per problem, where smaller $x$ means greater provisional difficulty. A rank label is an anchor at that problem's section: \verb|\ref{rank:x}| returns the section number, not the literal difficulty number. The visible \texttt{rank:x} text and the index give the difficulty number itself.

The entries are \emph{parameterized mathematical research schemes}, including empirical prediction, causal identification, historical reconstruction, theory and engineering. They are not 330 certified unsolved mathematical conjectures. A rigorous instance fixes the objects and restrictions below; changing those objects changes the question. Established identities, finite calculations, non-identification constructions and solved special cases are not being posed anew as open. Some residual tasks concern informative bounds, biological realization or external validation rather than a new theorem.

Every main definition and problem statement received a mathematical editorial review. Each entry now has an additional binding-conditions paragraph; direct corrections are separately recorded in the accompanying audit. The alternative formulations inherit the applicable conventions, but are different targets unless equivalence is proved. Where an old broad phrase conflicts with an explicit restriction below it, the explicit restriction controls.

\textbf{Literature boundary.} Thirteen selective primary-literature checks cover eighteen entries in this revision; their access levels and scope are recorded explicitly. The supplied bibliography and its earlier source-access descriptions are retained as provenance, not represented as newly verified. Neither this review nor a source's description proves the absence of later solutions. The date of this edition is 13 October 2026; there is no claim of exhaustive current-status certification for all 330 entries.

\subsection*{Binding conventions for a concrete instance}
\paragraph{Objects, quantifiers and non-vacuity.}
An instance declares measurable state, input, output and parameter spaces; a fixed candidate family; admissible initial laws, interventions and environments; a finite horizon or specified asymptotic regime; observation and sampling laws; and performance metrics, positive normalization scales and positive tolerances. These are chosen before fitting or testing, not adjusted after seeing failures. Shared parameters must genuinely be shared across the declared validation domain. Finite dimension alone does not restrict arbitrary real-number encodings: specify a finite library, bounded description length, or regularity/complexity bounds. Exclude constant or no-signal solutions through a fixed informative-output or baseline-improvement requirement where appropriate. A universally quantified result cannot be established by choosing one trivial favorable instance.

\paragraph{Probability laws and data.}
Every expectation or probability refers to the specified stochastic law, including initial uncertainty, process noise, observation noise and sampling as appropriate. A deterministic equation is not a probability model without a declared random input or initial law. ``Observed'' targets in population-level formulas mean population observables $g_j^*$, not noisy estimates asserted to be exact. Let $\widetilde g_j$ be independent validation estimates and suppose
\[
\mathbb{P}\{\lvert\widetilde g_j-g_j^*\rvert\le r_j\text{ for every }j\}\ge1-\alpha.
\]
Then an uncertainty-aware sufficient certification rule for positive target tolerances $\varepsilon_j$ is
\[
\max_j\frac{|\widehat g_j-\widetilde g_j|+r_j}{\varepsilon_j}\le1.
\]
On that simultaneous confidence event it implies every population error is at most its tolerance, by the triangle inequality. Vector/norm versions use corresponding valid simultaneous bounds. Training uncertainty must also be included when the claim averages over repeated training data. Evaluation on finitely many conditions does not establish a continuous-domain supremum without covering and regularity bounds.

\paragraph{Dynamics and well-posedness.}
For ODEs and SDEs give coefficient domains, initial laws and existence/uniqueness conditions in the intended solution sense, or keep the entire admitted solution set. For stochastic jump models specify nonnegative off-diagonal rates, diagonal normalization, rate orientation and non-explosion. For PDEs specify spatial domains, boundary/initial conditions, conservation or source terms, and a weak or classical solution class. Moving-domain concentration equations include dilution/transport terms. A stationary law is used only for a time-homogeneous augmented model under existence assumptions; a stationary approximation to a driven system must be justified. Absorbing extinction states are retained. Survival-conditioned, quasi-stationary and unconditional laws are distinct and must not be substituted for one another.

\paragraph{Identification and causal meaning.}
Two parameter values are observationally equivalent when they induce the same full distribution under all allowed observation/intervention designs, not merely the same likelihood value for one sample. A target is identified only if constant on the equivalence class. Gauge symmetries, label switching and reference conventions are quotiented or fixed. Finite-data compatible sets use a declared confidence, likelihood or discrepancy rule; an empty set signals model rejection. Bayesian credible sets are not automatically confidence sets. A prior or regularizer does not create data identification. Bounds intended to resolve a question must be informative at a predeclared width or decision threshold.

A causal contrast uses well-defined intervention versions and a stated causal model. Random assignment, or explicit consistency, overlap and exchangeability assumptions, is required for causal interpretation from observations. Interference and time-varying confounding are represented when present. Fixing a mediator by intervention differs from conditioning on its observed value. Marginal potential-outcome laws do not determine individual cross-world effects or an arbitrary joint coupling. Ethically admissible designs, rather than unconstrained interventions, define the experimental menu.

\paragraph{Metrics, units, optimization and limits.}
Specify the sampling or reference measure for densities and integrals, the log base, and dimensionless arguments of logarithms. Total variation is $\sup_A|P(A)-Q(A)|$; Wasserstein distances require the declared ground metric and appropriate moment conditions; KL divergence uses its stated absolute-continuity convention. All normed readouts are scaled consistently, and error probabilities use $0<\alpha<1$. A ratio has a positive denominator or an explicit undefined/zero convention. Thermodynamic entropy production is nonnegative in expectation under the relevant hypotheses; individual stochastic realizations need not be. An infimum is not an attained minimum without compactness/coercivity and continuity or another existence argument. Fix the order of large-population, long-time and vanishing-mutation limits. Stability claims declare perturbation space and symmetries.

\subsection*{How the provisional difficulty order was assigned}
Four editorial components are scored from 1 (lower burden) to 5 (higher burden): (a) formal or inferential difficulty, (b) access to evidence and identification burden, (c) mechanistic/computational integration, and (d) independent validation and transfer. Their unweighted total orders the entries. For selected entries in the highest bands, an explicit editorial priority breaks equal totals; all other ties are broken by original problem number. Equal totals are \emph{not} evidence of reliably different true difficulties. Scores are judgments, not measurements, completion-time estimates or consensus rankings.

The object ranked is an informative, nondegenerate resolution of the \emph{scoped} scientific target, with the stated budget and validation domain. A sharp useful impossibility result can qualify; an already known identity, vacuous bound, arbitrary favorable model or lookup table cannot. Rank 330 is merely the least difficult relative assessment in this catalogue, not a claim of an easy biological solution. Broad historical titles do not automatically outrank precisely formulated mechanistic tasks.



\end{document}
